\begin{helper}[Windowed short-edge count transported along an affine reindexing of breakpoints]
  Let $k,n,p,q,c,d \in \mathbb{N}$, let $\delta,\sigma,a',b' \in \mathbb{R}$, and let
  $s, s_A \colon \mathbb{N} \to \mathbb{R}$. (Throughout, $i-1$, $p-1$ and $i+c-d$ denote
  truncated natural-number subtraction.) Assume:
  \begin{itemize}
    \item $1 \le p$ and $d + 1 \le p + c$;
    \item (affine correspondence) for every $i \in \mathbb{N}$ with $p-1 \le i \le q$ we have
      $i + c - d \le k$ and
      \[
        s_A(i) = s\bigl(i + c - d\bigr) + \sigma
        ;
      \]
    \item (window count)
      \[
        \#\set{ j \in \set{1,\dots,k} \colon a'-\sigma \le s(j-1),\ s(j) \le b'-\sigma,\
          s(j)-s(j-1) < \delta }
        \le n
        .
      \]
  \end{itemize}
  Then
  \[
    \#\set{ i \in \set{p,\dots,q} \colon a' \le s_A(i-1),\ s_A(i) \le b',\
      s_A(i) - s_A(i-1) < \delta } \le n
    .
  \]

  \paragraph{Role.} This is the reusable form of the pull-back step used in
  \noderef{ArcShortWindowEdgeCount} and \noderef{ArcDeltaEpsNCount}: when the breakpoints of a
  circular arc's resolution are, on some index stretch, the breakpoints of the ambient curve's
  resolution under a single index map $i \mapsto i+c-d$ and a single value shift $\sigma$, the
  arc's short edges lying in a window $[a',b']$ are literally short edges of the ambient curve
  lying in the corresponding shifted window, so the ambient curve's own counting clause bounds
  them. The hypotheses mention only the two breakpoint sequences, so the same statement covers the
  ordinary (non-wrapping) case and both degenerate ends of the wrapping case, which differ only in
  the values of $(p,q,c,d,\sigma)$: the index map specialises to $i \mapsto j_0+(i-1)$ at
  $(c,d)=(j_0,1)$ and to $i \mapsto i-m$ at $(c,d)=(0,m)$. The window $[a',b']$ and the index range
  $\set{p,\dots,q}$ are free parameters rather than the extreme choices
  $[s_A(1),s_A(q)]$, $p=2$, which is what lets a consumer count inside a \emph{small} sub-window of
  a long breakpoint sequence -- the situation of \noderef{ArcDeltaEpsNCount}, where the window
  length is bounded by $2\epsilon^{-1}\delta$ but the whole interior window's length is not.
\end{helper}
\begin{proof}
  Put $\varphi(i) := i+c-d$ and let
  \[
    S := \set{ i \in \set{p,\dots,q} \colon a' \le s_A(i-1),\ s_A(i) \le b',\
      s_A(i) - s_A(i-1) < \delta },
  \]
  \[
    T := \set{ j \in \set{1,\dots,k} \colon a'-\sigma \le s(j-1),\ s(j) \le b'-\sigma,\
      s(j)-s(j-1) < \delta }
    .
  \]
  We show $\varphi$ maps $S$ into $T$ and is injective on $S$; then $\#S \le \#T \le n$, which is
  the claim.

  \emph{No truncation on the range.} Let $i$ with $p \le i \le q$. Since $p \ge 1$ we have
  $i \ge 1$, so $i-1$ is ordinary subtraction and $(i-1)+c = (i+c)-1$. Since $i \ge p$ we have
  $i + c \ge p + c \ge d+1$, so $\varphi(i) = i+c-d \ge 1$ is ordinary subtraction, and
  $(i-1)+c-d = (i+c)-1-d = \varphi(i)-1$, again ordinary subtraction.

  \emph{$\varphi$ maps $S$ into $T$.} Let $i \in S$, so $p \le i \le q$, $a' \le s_A(i-1)$,
  $s_A(i) \le b'$ and $s_A(i)-s_A(i-1) < \delta$. Both $i$ and $i-1$ lie in $\set{p-1,\dots,q}$
  (for $i$: $p-1 \le p \le i \le q$; for $i-1$: $i-1 \ge p-1$ since $i \ge p$, and
  $i-1 \le i \le q$), so the affine correspondence applies at both indices and gives
  \[
    \varphi(i) = i+c-d \le k,
    \qquad
    s\bigl(\varphi(i)\bigr) = s_A(i) - \sigma,
    \qquad
    s\bigl((i-1)+c-d\bigr) = s_A(i-1) - \sigma
    .
  \]
  By the no-truncation paragraph, $(i-1)+c-d = \varphi(i)-1$, so the last identity reads
  $s(\varphi(i)-1) = s_A(i-1)-\sigma$; and $\varphi(i) \ge 1$, so together with $\varphi(i) \le k$
  we get $\varphi(i) \in \set{1,\dots,k}$.

  It remains to verify the three defining conditions of $T$ at $j = \varphi(i)$.
  \begin{itemize}
    \item $a'-\sigma \le s(\varphi(i)-1)$: indeed $s(\varphi(i)-1) = s_A(i-1)-\sigma \ge
      a'-\sigma$, using $a' \le s_A(i-1)$.
    \item $s(\varphi(i)) \le b'-\sigma$: indeed $s(\varphi(i)) = s_A(i)-\sigma \le b'-\sigma$,
      using $s_A(i) \le b'$.
    \item $s(\varphi(i)) - s(\varphi(i)-1) = \bigl(s_A(i)-\sigma\bigr) - \bigl(s_A(i-1)-\sigma\bigr)
      = s_A(i) - s_A(i-1) < \delta$, since $i \in S$.
  \end{itemize}

  \emph{$\varphi$ is injective on $S$.} If $i_1, i_2 \in S$ with $\varphi(i_1) = \varphi(i_2)$,
  then by the no-truncation paragraph both $\varphi(i_1) = i_1+c-d$ and $\varphi(i_2) = i_2+c-d$
  are ordinary subtractions of $d$ from $i_1+c$ resp.\ $i_2+c$, so $i_1+c = i_2+c$ and hence
  $i_1 = i_2$.

  Hence $\#S \le \#T \le n$.
\end{proof}
